\chapter*{Résumé}
\addcontentsline{toc}{chapter}{Résumé}

Dans un environnement organisationnel marqué par la diversification des compétences, l'évolution rapide des projets et la nécessité d'un pilotage transversal, les entreprises ont besoin d'outils capables de centraliser l'information et d'anticiper les déséquilibres. Le présent travail s'inscrit dans ce contexte et porte sur Strategic Copilot : plateforme SaaS intelligente d'aide à la décision pour la gestion stratégique des talents et des projets. La plateforme accompagne les managers, les responsables RH et les talents au moyen de tableaux de bord interactifs, d'alertes intelligentes et de recommandations issues de l'intelligence artificielle, afin de piloter efficacement les compétences, les charges de travail et le portefeuille de projets.

L'objectif de ce projet de fin d'études est de concevoir et de réaliser une solution web intégrée offrant la visualisation des indicateurs RH et projet, la détection des risques (surcharge, fin de contrat, manque de compétences), l'analyse de la viabilité des projets, la génération de recommandations automatisées par IA et la gestion des demandes RH au sein de workflows décisionnels automatisés, avec une expérience adaptée à chaque profil utilisateur.

La démarche adoptée combine l'analyse des besoins, la définition des cas d'usage par espace métier (manager, RH, talent), la conception de l'architecture applicative, puis le développement itératif d'une interface React.js avec Vite et Tailwind CSS. L'intégration repose sur Supabase et PostgreSQL pour la persistance, l'authentification et les API, sur n8n pour l'orchestration des flux, et sur le modèle Llama pour l'assistance décisionnelle. Des phases de tests et de validation ont permis d'ajuster les fonctionnalités sur les scénarios prioritaires.

Les travaux ont abouti à une plateforme opérationnelle structurée en espaces dédiés : pilotage des projets et des risques pour les managers, gestion des talents, des écarts de compétences et des demandes pour les RH, suivi des missions, tâches et formations pour les talents. Les tableaux de bord, les alertes et le copilot intégré confirment la pertinence de l'aide à la décision ; les workflows automatisés renforcent la traçabilité des arbitrages. L'ensemble démontre la faisabilité d'une solution SaaS au service du pilotage stratégique des ressources humaines et des projets.

\textbf{Mots-clés :} SaaS, Intelligence artificielle, Gestion des talents, Gestion de projets, Aide à la décision.

\clearpage
